\section{Conclusion}

We extract a common connectivity and weight-sharing constraint by
aligning source solutions, jointly projecting their task-specific
weights and selecting a finite bank by prediction error after numerical
refitting. New tasks receive only the selected constraint; their
numerical parameters are learned afresh.

For aligned source weights with fixed active connections, we characterize
exact recovery through scale-invariant magnitudes and sign profiles.
A condition on the graph of co-occurring codes gives uniqueness even
when no pair of codes repeats across neurons.
For partial observations, a completion criterion and a functional
ReLU example show how several individually ambiguous tasks can
determine the groups jointly. Recovery from erroneous pairwise
relations additionally requires sufficient redundancy between
the observed parts.

Joint projection and direct pair recovery both identify the full signed
structure in the two clean synthetic geometries. On the class-disjoint
image tasks, both use 114 numerical parameters per new task instead of
12,594 for Dense, with lower prediction accuracy at the tested code
budgets. Direct pair recovery also trails G34 on both datasets.
Both select the largest tested codebook, so these results do not
establish an optimal sharing budget.
Good source prediction alone therefore does not establish a useful
sharing constraint for new classes. Candidate coverage, robustness to
weight estimation error and the appropriate strength of sharing remain
central practical questions.
